\clearpage
\SourceLeaf{315}
\GaussVSourceUnnumbered
\pdfbookmark[1]{Göttingische gelehrte Anzeigen. 1841 Januar 18.}{anzeigen-1841-januar-18}
\begin{center}
\rule{.70\linewidth}{.4pt}\\[.7\baselineskip]
Göttingische gelehrte Anzeigen. 1841 Januar 18.\\[.7\baselineskip]
\rule{.70\linewidth}{.4pt}
\end{center}

Am 10. December v. J. ist der Königl. Societät von dem Hofrath \textsc{Gauss}
eine Vorlesung mit der Ueberschrift:
\begin{center}\textit{Dioptrische Untersuchungen}\end{center}
\noindent
überreicht, von welcher hier ein kurzer Bericht abzustatten ist.

Die Betrachtung des Weges, welchen durch eine beliebige Anzahl auf einer
gemeinschaftlichen Axe geordneter Linsengläser solche Lichtstrahlen nehmen, die
gegen die Axe sehr wenig geneigt sind, und der davon abhängenden Erscheinungen,
bietet viele durch Allgemeinheit und Eleganz merkwürdige Resultate
dar, mit deren Aufsuchung sich mehrere der ersten Mathematiker beschäftigt
haben. Der von \textsc{Cotes} aufgefundene Lehrsatz wurde seiner Zeit mit einer Art
von Bewunderung aufgenommen, und man hat sogar auf denselben (obwohl
vielleicht nur durch ein Missverständniss der darauf bezüglichen Stelle in \textsc{Smith}’s
Optik) den bekannten ehrenden Nachruf bezogen, mit welchem \textsc{Newton} den
frühen Tod jenes genialen Geometers beklagte. \textsc{Euler}’s Arbeiten umfassten alle
Theile der Dioptrik. Ganz besonders aber ist hierher zu rechnen die schöne
und prägnante Behandlungsart des in Rede stehenden Gegenstandes durch \textsc{Lagrange}
in den Schriften der Berliner Academie der Wissenschaften von 1778,
wozu später noch schätzbare Zusätze theils von ihm selbst, theils von \textsc{Piola}
und \textsc{Möbius} hinzu gekommen sind. Nach solchen Arbeiten könnte die Ernte
auf diesem Felde vollendet
\clearpage
\SourceLeaf{316}
\GaussVSourceFolio{310}{L}{ANZEIGEN.}
\noindent
scheinen: gleichwohl bleibt noch Mehreres zu wünschen übrig. Allen
jenen Lehrsätzen liegt die Voraussetzung zum Grunde, dass die Dicke der Linsen
als unendlich klein betrachtet werde, eine Voraussetzung, an welche man so
gewöhnt ist, dass sie meistens, als sich von selbst verstehend, gar nicht einmal
erwähnt wird, mit welcher aber bei der Anwendung auf wirkliche Fälle jene
Lehrsätze nur Annäherungen, zuweilen nur rohe Annäherungen bleiben. Diese
fast allgemeine Vernachlässigung der Dicke der Linsen erstreckt sich zum Theil
selbst auf die ersten Begriffsbestimmungen der Dioptrik, bei welchen der mathematische
Sinn von einer schwankenden Unbestimmtheit unangenehm berührt
wird. Wenn von der Brennweite einer Linse ohne nähere Bestimmung geredet
wird, so weiss man noch nicht, ob die Entfernung des Brennpunktes von der
nächsten Oberfläche der Linse, oder vom sogenannten optischen Mittelpunkte
derselben, oder von dem zwischen beiden Oberflächen in der Mitte liegenden
Punkte, oder die von allen diesen Bestimmungen wieder verschiedene Grösse
gemeint ist, welche bei der Vergleichung der scheinbaren Grösse eines unendlich
entfernten Gegenstandes mit seinem Bilde zum Grunde gelegt werden muss.

Es ist nun zwar nicht zu leugnen, dass in sehr vielen Fällen die Vernachlässigung
der Dicke der Linsen zulässig, oder nützlich, oder sogar nothwendig
sein kann, wo entweder an grösserer Schärfe nichts gelegen ist, oder, wo eine
die Dicke der Linsen mit berücksichtigende Schärfe in unerträgliche Weitläufigkeiten
führen würde, oder wo eine scharfe Rechnung durch vorgängige genäherte
Überschläge erst vorbereitet werden soll: allein eben so gewiss ist, dass die
Würde der Wissenschaft Präcision in ihren Begriffsbestimmungen erfordert, und
dass man sich der Aufopferung der Schärfe gern in allen Fällen enthoben sieht,
wo es ohne erheblichen oder ohne allen Nachtheil für die Einfachheit und
Geschmeidigkeit der Resultate möglich ist.

Indem die vorliegende Abhandlung zum Hauptzwecke haben sollte, zu zeigen,
dass den oben angedeuteten Lehrsätzen eine Erweiterung gegeben werden kann,
in der sie unter strenger Berücksichtigung der Dicke der Linsen gar nichts von
ihrer Einfachheit verlieren, war es nothwendig, auf die ersten Grundlehren der
Dioptrik in einer neuen Darstellungsart zurück zu kommen. Dieser, und dem
ganzen Inhalte der Abhandlung Schritt vor Schritt zu folgen, kann hier nicht
der Ort sein, wäre auch um so überflüssiger, da die Abhandlung selbst bereits
gedruckt ist. Wir beschränken uns daher hier auf den Versuch, anschaulich zu
\clearpage
\SourceLeaf{317}
\GaussVSourceFolio{311}{R}{DIOPTRISCHE UNTERSUCHUNGEN.}
\noindent
machen, wie der Fall eines Systems von Linsen mit endlicher Dicke auf den
Fall eines Systems unendlich dünner Linsen zurück geführt wird.

An dem Wege jedes Lichtstrahls, der durch eine Glaslinse geht, ohne mit
ihrer Axe zusammen zu fallen, sind drei Theile zu unterscheiden; der erste, vor
dem Eintritte in die Glaslinse; der zweite, innerhalb derselben; der dritte, nach
dem Austritte. In jedem Systeme von Lichtstrahlen, die unter sich parallel,
oder aus Einem Punkte divergirend, oder nach Einem Punkte zu convergirend,
auf eine Glaslinse fallen, ist Einer, von dessen Wege der dritte Theil dem
ersten parallel wird: ein solcher Strahl heisst ein Hauptstrahl. Sämmtliche
Hauptstrahlen haben die merkwürdige Eigenschaft, dass der zweite Theil ihres
Weges, nöthigenfalls vorwärts oder rückwärts geradlinig verlängert, die Axe der
Linse in einem und demselben Punkte trifft: diesen Punkt haben einige Schriftsteller
über Dioptrik den optischen Mittelpunkt der Linse genannt. Er liegt
innerhalb der Linse, wenn sie convex-convex oder concav-concav ist; in der
krummen Oberfläche, wenn sie plan-convex oder plan-concav ist; ausserhalb, und
zwar immer auf der Seite der stärkern Krümmung, wenn die Linse convex-concav
oder concav-convex ist. Diese Eigenschaft macht zwar den erwähnten Punkt
allerdings merkwürdig; allein da er sonst gar keine practisch nützliche Brauchbarkeit
hat, so ist die Belegung desselben mit einer besondern Benennung eine
kaum verdiente Auszeichnung, die vielleicht sogar dadurch nachtheilig geworden
ist, dass sie, wie es scheint, hie und da zu dem Irrthume verleitet hat, als ob
man die bekannten einfachen Relationen, welche zwischen einem Objecte und
seinem Bilde für eine Linse von unendlich kleiner Dicke Statt finden, auf eine
Linse von endlicher Dicke ohne weiteres bloss vermittelst der Beziehung auf
ihren optischen Mittelpunkt übertragen, oder mit anderen Worten, für eine Linse
von endlicher Dicke eine andere von unendlich kleiner Dicke in dem optischen
Mittelpunkte der erstern substituiren dürfte.

Eine ganz andere Wichtigkeit haben hingegen diejenigen zwei Punkte, wo
der erste und der dritte Theil des Weges eines Hauptstrahls, nöthigenfalls vorwärts
und rückwärts verlängert, die Axe der Linse schneiden. Auch sie haben
die Eigenschaft, für sämmtliche Hauptstrahlen dieselben zu sein, und verdienen
durch besondere Benennungen ausgezeichnet zu werden: der Verfasser nennt sie
den \textit{ersten} und den \textit{zweiten Hauptpunkt der Linse}. Sie haben aber zugleich die
wichtige und wie es scheint bisher nicht bemerkte Eigenschaft, dass alle ausfahrenden
\clearpage
\SourceLeaf{318}
\GaussVSourceFolio{312}{L}{ANZEIGEN.}
\noindent
Strahlen sich relativ gegen den zweiten Hauptpunkt genau so verhalten,
wie sie sich gegen den ersten verhalten würden, wenn die einfallenden Strahlen
anstatt der wirklichen Linse eine andere von unendlich kleiner Dicke und von
derselben Brennweite im ersten Hauptpunkte befindliche träfen. Als Brennweite
der wirklichen Linse gilt hier die Entfernung ihres Brennpunktes, d. i. des Vereinigungspunktes
der parallel mit der Axe einfallenden Strahlen, von dem zweiten
Hauptpunkte, und eben diese Grösse ist es auch, die bei der Vergleichung der
scheinbaren Grösse eines unendlich entfernten Gegenstandes mit der Grösse
seines Bildes zum Grunde gelegt werden muss. Hiedurch erhält also der Begriff
der Brennweite eine scharf bestimmte Haltung, und es mag zugleich bemerkt
werden, dass dann die Brennweite \textit{dieselbe} bleibt, die Strahlen mögen von der
einen oder von der andern Seite einfallen, nur dass natürlich im zweiten Falle
der vorige erste Hauptpunkt an die Stelle des zweiten tritt. Man darf also,
so lange man sich auf die gegen die Axe sehr wenig geneigten Strahlen
beschränkt (oder von der Abweichung wegen der Kugelgestalt abstrahirt), alle
Rechnungen über Linsen von endlicher Dicke ganz eben so führen, als wäre
die Dicke unendlich klein und der Zwischenraum zwischen den beiden Hauptpunkten,
in deren einem man sich die ideale Linse vorstellt, gleichsam vernichtet.
Übrigens ist dieser Zwischenraum nahe dem dritten Theile der Dicke der Linse
gleich, wenn sie von gewöhnlichem Glase, und etwas grösser (nahe \(\frac{5}{13}\) der
Dicke), wenn sie von Flintglas ist.

Die Begriffe von Hauptstrahlen und Hauptpunkten und deren Anwendungen
lassen sich auch auf ein System von mehreren Linsen auf gemeinschaftlicher
Axe ausdehnen, während von einem optischen Mittelpunkte in der obigen Bedeutung
dann gar nicht mehr die Rede sein kann. Für ein achromatisches Objectiv,
dessen beide Bestandtheile einander beinahe berühren, als ein Ganzes betrachtet,
wird der Abstand der beiden Hauptpunkte von einander sehr nahe der Summe
der beiden respectiven Abstände in den einzelnen Linsen gleich.

Über die Methoden, welche in der Abhandlung zur Bestimmung der Brennweiten
von Linsengläsern obigen Grundlagen gemäss entwickelt werden, können
wir uns hier nicht verbreiten. Am Schlusse der Abhandlung sind noch Bemerkungen
beigefügt, wodurch die eigenthümlichen Vorzüge, deren die sogenannten
dialytischen Fernröhre in Beziehung auf Farbenreinheit fähig sind, in ihr wahres
Licht gesetzt werden.
\begin{center}\rule{.26\linewidth}{.4pt}\end{center}
